% Einheit 4 von 4 – Abstand als Argument (bank/abstaende/e4.jsonl, zusammenbau v0.1)
%% TODO Zeitmarke Einheit 4: Marken-Zeile nicht lesbar
%% TODO Zweigzeile Teil 1: was hier gelernt wird (Satz in Schülersprache aus dem Einheitstitel) – Einheit 4
\einheitenkopf[e4]{Einheit 4 von 4 · Abstand als Argument}
\zweigzeile{Abitur GK}
\setcounter{aufgabe}{15}

%% TODO Titel: Ich-kann-Satz fehlt in der Bank (Platzhalter: Kettenname) – Nr. 16
%% TODO Anweisung: ein Satz über der ersten Teilaufgabe fehlt in der Bank – Nr. 16
\begin{aufgabe}{Abstand als Argument}
%% TODO \rechenplatz{n}: Schrittzahl des Grundfalls fehlt in der Bank (nur bei fester Schreibform, 2.3 b)
\begin{teile}
\teil Senkrecht oder schräg? Die Strecke verbindet P(1 | 2 | 4) mit dem Punkt Q(1 | 2 | 0) der xy-Ebene. Kreuze an, rechne nicht. \\ \kreuz{senkrecht: die Strecke ist der Abstand} \\ \kreuz{schräg: die Strecke ist länger als der Abstand}
\teil B(1 | 0 | 4) liegt in der Ebene $L_1$: $x + 2y + 2z = 9$, E(1 | 0 | 1) in der parallelen Ebene $L_2$: $x + 2y + 2z = 3$. Ist der Abstand von B und E kleiner, größer oder genauso groß wie der Abstand der Ebenen? Begründe. \hfill (Abitur 2023 GK)
\teil Die Abbildung zeigt eine Tribüne im Querschnitt als Strecke GJ auf der Geraden t und den Punkt F. F hat von t etwa den Abstand 3,1. Begründe mit einer Zeichnung, dass kein Punkt der Strecke GJ von F den Abstand 3,1 hat. \hfill (Abitur 2018 LK)

\begin{ksys}[xmin=0,xmax=8,ymin=0,ymax=8]
\gerade{0.5}{0}{t}
\punkt{4}{2}{G}
\punkt{8}{4}{J}
\punkt{1}{4}{F}
\end{ksys}
\teil Eine senkrecht stehende rechteckige Werbetafel ABED ist an einem Mast befestigt, dessen Achse die z-Achse ist. Die oberen Ecken D und E liegen symmetrisch zu einer senkrechten Ebene durch die z-Achse, M ist der Mittelpunkt der Oberkante DE. Begründe, dass der Abstand der Tafel zur Mastachse gleich dem Abstand von M zur Mastachse ist. \hfill (Abitur 2023 LK)
\teil Eine Zeltwand liegt in der Ebene W: $3y + 2z = 15$, ihre Bodenkante in der xy-Ebene bei $y = 5$; 1 LE = 1 m. Aus der Wand wird ein Rechteck von der Bodenkante bis zur Höhe 1,20 m ausgeschnitten und zu einem waagerechten Vordach nach außen geklappt. Wie lang ist das Vordach? \hfill (Abitur 2017 LK) \\ \leerfeld[m]
\teil Ein Zelt hat die Form einer Pyramide mit quadratischem Boden und der Spitze S(5 | 3 | 4) senkrecht über dem Bodenmittelpunkt M(5 | 3 | 0); eine Wand liegt in der Ebene W: $4y - 3z = 0$; 1 LE = 1 m. Eine Lampe soll von allen vier Wänden 1 m entfernt sein. Wo hängt sie? \hfill (Abitur 2017 LK) \\ L( \leerfeld | \leerfeld | \leerfeld )
\teil Der Punkt P(0 | 0 | p) im Innern einer Pyramide soll von allen Seitenflächen und von der Grundfläche (xy-Ebene) gleich weit entfernt sein; eine Seitenfläche liegt in F: $2y + z = 4$. Vorgelegt: I $\overrightarrow{OQ} = (0 | 0 | p) + t \cdot (0 | 2 | 1)$; II $2 \cdot 2t + 1 \cdot (p + t) = 4$; III $|\overrightarrow{PQ}| = p$. Erläutere die geometrischen Überlegungen. \hfill (Abitur 2024 LK)
\end{teile}
\end{aufgabe}

%% TODO Titel: Ich-kann-Satz fehlt in der Bank (Platzhalter: Kettenname) – Nr. 17
%% TODO Anweisung: ein Satz über der ersten Teilaufgabe fehlt in der Bank – Nr. 17
\begin{aufgabe}{Abstand als Argument – Fehler finden, Begründen, Anwendung}
\begin{teile}
\teil Finde den Fehler. A(0 | 1 | 5) liegt in $L_1$: $2x + y + 2z = 11$, B(1 | 0 | 0) in der parallelen Ebene $L_2$: $2x + y + 2z = 2$. Jan schreibt: \rechnung{\text{Abstand der Ebenen} &= |\overrightarrow{AB}| = \sqrt{27} \approx 5{,}20}
\teil Warum ist das Lot die kürzeste Verbindung von einem Punkt zu einer Ebene? Begründe.
\teil Auf einem Spielplatz stehen zwei senkrechte Pfosten mit den Fußpunkten A(2 | 1 | 0) und B(8 | 4 | 0); 1 LE = 1 m. Eine Laterne soll auf der Linie durch A und B stehen und vom Pfosten in A doppelt so weit entfernt sein wie vom Pfosten in B. Wo kann sie stehen? ( \leerfeld | \leerfeld | \leerfeld ), ( \leerfeld | \leerfeld | \leerfeld )
\end{teile}
\end{aufgabe}
